\section{Study 2: the allocation causes the later decision}\label{sec:causal}

Study 1 shows where the budget goes. It does not show that this matters. The
natural test --- let the situation shift after the budget is spent, then see
whether episodes that recovered the caveat behave differently --- is
observational: verification is chosen by the model, so the two groups differ in
whatever made the model choose.

An earlier two-phase experiment ran exactly that comparison and reported a large
split. Auditing its implementation we found two paths that the design left open
and the reporting did not separate. First, the triage arm's system prompt was
constructed once and reused for the later call, so the randomised instruction
was still in context when the second decision was made and could act on it
directly. Second, the phase-1 conversation was carried forward, so the model saw
its own earlier reasoning --- which in the drifted arm frequently argued
explicitly against the corrupted direction. ``Verified'' and ``unverified''
episodes therefore differed in what the model had already committed to in
writing, not only in what evidence it held.

We close both by construction, pre-registering two experiments before any call.

\subsection{A randomised-instrument replay}

For each of the \ThreeAN{} episodes of the earlier experiment we replay
\emph{only} the second decision, as a fresh session: the triage sentence is
absent from every arm, no conversation history is carried, the memory block and
its per-episode order are reproduced, and the only thing that differs is a
\textsc{recovered} block holding source records for exactly the memories that
episode's budget was actually spent on. The situation text is byte-identical
across arms.

Triage assignment was randomised upstream and moves recovery hard. In the replay
it cannot reach the outcome by any route except through what was carried
forward, because the session is otherwise identical across arms. This is a
randomised encouragement design in which the exclusion restriction is true by
construction rather than assumed.

It replicates. The corrupted-direction rate is \ThreeATrK/\ThreeATrN{}
(\ThreeATrPct\%) under triage against \ThreeADrK/\ThreeADrN{} (\ThreeADrPct\%)
without it (CMH stratified by model, $p = \ThreeACMH$). And the exclusion
restriction can be checked rather than argued: matched on what phase~1
recovered, the two arms are indistinguishable in both strata ---
\ThreeARecDrK/\ThreeARecDrN{} versus \ThreeARecTrK/\ThreeARecTrN{} among
episodes that recovered the caveat ($p = \ThreeARecCMH$), and
\ThreeAUnrecDrK/\ThreeAUnrecDrN{} versus \ThreeAUnrecTrK/\ThreeAUnrecTrN{} among
those that did not ($p = \ThreeAUnrecCMH$). With the instruction removed from
the room, arms holding the same evidence behave the same way.

The Wald estimate of the effect of recovering the corrupted memory's source is
$\IVEst$ (95\% CI $[\IVLo, \IVHi]$ by stratified bootstrap over models): close
to complete compliance, an agent that recovers the caveat essentially does not
walk into the corrupted direction.

\paragraph{The original design was conservative, not inflated.} Removing the
carried-forward conversation \emph{raises} the corrupted-direction rate among
episodes that never recovered the caveat, from 79\% in the original to
\ThreeAUnrecDrPctText\% here. The model's own earlier prose had been suppressing
the choice, and the original design credited that suppression to provenance. The
confound worked against the reported effect.

\subsection{Randomising the provenance itself}

The replay still lets the model choose what it recovered. So we also assign it.
In \ThreeBN{} fresh single-call episodes the memory block is the drifted lineage
in both arms, and two source records are carried forward: in \textsc{carry-target}
one of them is the corrupted memory's, in \textsc{carry-other} neither is. Both
arms carry exactly two records; only their identity differs, and it is assigned
by seeded PRNG.

\begin{table}[t]
\centering\small
\caption{Study 2b. Corrupted-direction rate when the provenance carried into the
later decision is randomised rather than chosen. Both arms carry two source
records; only whether one of them is the corrupted memory's differs.}
\label{tab:carry}
\begin{tabular}{lrr}
\toprule
& corrupted memory's source carried & not carried \\
\midrule
Opus 5 & 0/25 & 14/25 \\
Sonnet 5 & 2/25 & 25/25 \\
Haiku 4.5 & 0/25 & 25/25 \\
Sol & 0/25 & 25/25 \\
Terra & 1/25 & 25/25 \\
Luna & 0/25 & 25/25 \\
\midrule
pooled & 3/150 (2\%) & 139/150 (93\%) \\
\bottomrule
\end{tabular}
\end{table}


Table~\ref{tab:carry} gives the result: \ThreeBTgtK/\ThreeBTgtN{} against
\ThreeBOthK/\ThreeBOthN{}, a difference of \ThreeBDiff{} percentage points
(CMH $\chi^2 = \ThreeBChi$, $p \ThreeBCMH$), in the same direction in
\ThreeBAllModels{} of six models, with five of six choosing the corrupted
direction in every single episode where its provenance was absent.

\paragraph{Randomisation reaches severity, not only direction.} Unguarded moves
into the corrupted direction --- anything beyond a small guarded test --- occur
in \ThreeBNGOth{} episodes when the corrupted memory's provenance is withheld
and \ThreeBNGTgt{} when it is carried (CMH $\chi^2 = \ThreeBNGChi$,
$p = \ThreeBNGCMH$). Every unguarded choice in the experiment is on the side
that never saw the caveat. The five-clause strict harm rule still returns zero
because self-reported hedging saturates: \ThreeBNGHedged{} of those choices omit
the uncertainty flag.

Two independent identification strategies --- an instrument in 2a, full
randomisation in 2b --- put the effect at the same place, near $-\ThreeBDiff$ percentage
points. What the budget happened to recover is not a marker of some other
difference between episodes. It is the thing that decides.

\paragraph{What this does and does not establish.} It establishes that
\emph{holding} the recovered caveat changes the later decision. It does not
establish that models allocate budget \emph{in order to} secure future evidence;
Study 1 says the opposite, that they allocate toward the present plan. The
causal chain we have closed is allocation $\rightarrow$ what is in context later
$\rightarrow$ behaviour, not any claim about foresight.
